\documentclass[11pt]{article}
\usepackage{amsmath,amssymb,amsthm}
\usepackage[margin=1in]{geometry}
\title{Appendix: Assumptions, Notation, and Theorems Used}
\date{}
\begin{document}
\maketitle

\section{Assumptions}
(A1) $\Omega$ is smooth bounded or periodic.\\
(A2) $D_\kappa,D_\tau,D_\Sigma>0$.\\
(A3) $\phi(\Sigma)\in[0,1]$ for $\Sigma\ge 0$, smooth.\\
(A4) Initial data satisfy $\kappa_0\in[0,1]$, $\tau_0\in[0,1]$, $\Sigma_0\ge 0$.

\section{Theorems invoked}
Maximum principle for scalar parabolic PDE implies:
if reaction points inward at boundaries of an interval, the interval is forward invariant.

Comparison principle yields global $L^\infty$ bounds for $\Sigma$ via a linear supersolution.

Routh--Hurwitz criteria for cubic polynomials:
$\lambda^3+a_1\lambda^2+a_2\lambda+a_3$ has roots with negative real part iff
$a_1>0$, $a_2>0$, $a_3>0$, and $a_1a_2>a_3$.

\end{document}
